\section{现象层：有限语言、循环子语言与最小核结构}

\subsection{本体窗口与线性稳定语言}

在有限窗离散化中，引入本体窗口空间
$$
\Omega_m=\{0,1\}^m,\qquad \abs{\Omega_m}=2^m.
$$
\begin{definition}[线性稳定语言]\label{def:stable_language}
\MathTag 定义线性稳定语言
$$
X_m=\left\{w\in\{0,1\}^m:\forall i\in\{1,\dots,m-1\},\ (w_i,w_{i+1})\neq(1,1)\right\}.
$$
\end{definition}
该集合亦可视为黄金均值移位的长度 $m$ 允许词集合。\cite{scholarpedia_symbolic_dynamics}

\subsection{线性稳定语言计数定理}

\begin{theorem}\label{thm:Xm_count}
\MathTag
\AuditTag 台账编号：\texttt{L-lang-01}。
$$
\abs{X_m}=F_{m+2}.
$$
\end{theorem}

\begin{proof}
\MathTag 依赖：第 2 节 Fibonacci 数列 $(F_n)$ 的定义；定义 \ref{def:stable_language}。
设 $a_m=\abs{X_m}$。按首位 $w_1$ 分类：
\begin{itemize}
  \item 若 $w_1=0$，则余下 $m-1$ 位可为任意 $X_{m-1}$ 中元素，计数为 $a_{m-1}$。
  \item 若 $w_1=1$，则必须有 $w_2=0$，余下 $m-2$ 位可为任意 $X_{m-2}$ 中元素，计数为 $a_{m-2}$。
\end{itemize}
因此递推
$$
a_m=a_{m-1}+a_{m-2},
$$
并由初值 $a_1=2$、$a_2=3$ 得 $a_m=F_{m+2}$。证毕。
\end{proof}

\subsection{循环子语言与 Lucas 计数}

定义循环稳定语言
$$
C_m=\left\{w\in\{0,1\}^m:\forall i\in\{1,\dots,m\},\ (w_i,w_{i+1})\neq(1,1)\right\},
\qquad w_{m+1}=w_1.
$$

\begin{theorem}\label{prop:Cm_Lm}
\MathTag
\AuditTag 台账编号：\texttt{L-lang-02}。
$$
\abs{C_m}=L_m.
$$
\end{theorem}

\begin{proof}
\MathTag 依赖：定理 \ref{thm:Xm_count}；第 2 节 Lucas 数列 $(L_n)$ 的定义。
当 $m=1$ 时，$C_1=\{0\}$，故 $\abs{C_1}=1=L_1$；当 $m=2$ 时，循环约束等价于线性无相邻约束，$C_2=X_2$，故 $\abs{C_2}=3=L_2$。以下设 $m\ge 3$。令 $c_m=\abs{C_m}$。按 $w_1$ 分类：
\begin{itemize}
  \item 若 $w_1=0$，则其余 $m-1$ 位仅需满足线性无相邻约束，计数为 $\abs{X_{m-1}}=F_{m+1}$。
  \item 若 $w_1=1$，则 $w_2=0$ 且 $w_m=0$，中间 $m-3$ 位满足线性无相邻约束，计数为 $\abs{X_{m-3}}=F_{m-1}$。
\end{itemize}
故
$$
c_m=F_{m+1}+F_{m-1}=L_m.
$$
证毕。
\end{proof}

圆环上无相邻 $1$ 的二进制配置与 Lucas 数列一致的事实在组合计数资料中作为典型例题出现。\cite{cameron_notes_on_counting}

\subsection{边界态集合与分解}

定义边界态集合
$$
B_m=X_m\setminus C_m.
$$
当 $m\ge 2$ 时，由恒等式 $L_m=F_{m+1}+F_{m-1}$ 与 $\abs{X_m}=F_{m+2}$ 得
$$
\abs{B_m}=\abs{X_m}-\abs{C_m}=F_{m+2}-L_m=F_{m-2}.
$$
并且
$$
X_m=C_m\sqcup B_m.
$$
当 $m=1$ 时，$\abs{B_1}=\abs{X_1}-\abs{C_1}=2-1=1$，与上述递推一致性不矛盾，但不再用 $F_{m-2}$ 书写以避免负指标。
该分解在后续接口一致性与耦合项分析中用于刻画“首尾粘合的缺陷来源”。

\subsection{最小核实例（m=6）}

当 $m=6$ 时
$$
\abs{\Omega_6}=2^6=64,\qquad \abs{X_6}=F_8=21,
$$
并由上述定理得
$$
\abs{C_6}=L_6=18,\qquad \abs{B_6}=F_4=3,
\qquad X_6=C_6\sqcup B_6,\qquad 21=18\oplus 3.
$$
该最小核用于后续章节的枚举验证与结构可视化基准。
